\chapter{Fundamentals still asked}
\label{ch:fundamentals}

\section*{What this chapter is for}

A build round in this domain is still a build round. The questions that decide
it are frequently ordinary engineering questions, and a candidate who has spent
a month on retrieval and cannot say what happens when two people buy the last
item has prepared the wrong half.

This chapter is a reminder rather than a course. \judgement{If you are senior
in backend or platform work you know all of it; the value is in having answered
each one recently, out loud, because fluency decays and these are asked under
time pressure.}

\section*{The questions}

\begin{bankq}{Two people try to buy the last one. What happens?}
\qasked
A scenario, and one of the most reliably asked questions in any build round
that touches inventory, booking or capacity.

\qtested
Whether you reach for a lock, a transaction, or a conditional write, and
whether you know why.

\qstaff
The failure is read-then-write: both read one remaining, both decide it is
available, both write zero, and two orders exist against one item. A
transaction alone does not fix it \dash{} at the isolation level most databases
default to, both transactions read the same value and both commit happily.

Three real answers, and knowing which you are giving is the point. A
conditional update \dash{} decrement \emph{where} the quantity is still
sufficient \dash{} which makes the check and the write one atomic operation and
lets the caller detect losing by finding no rows affected. A version column,
where the write requires the version you read and fails if it moved. Or raising
the isolation level, which pushes the problem into the database and produces
serialisation failures the application still has to handle.

\judgement{The conditional update is the boring default and should be said
first. It needs no extra column, no retry loop in the common case, and it
degrades to a clean failure. The version column earns its place when the
conflict is over an aggregate rather than a counter.}

And the constraint underneath it: a check that the quantity cannot be negative,
so the invariant survives a bug in any of the three.

\qsenior
``Wrap it in a transaction'' \dash{} the most common answer and insufficient at
the default isolation level, which is the specific thing being probed.

\qcontested
Optimistic against pessimistic, which is a real fork decided by contention
rather than by preference.

\qdrill
Write the conditional update, as SQL, from memory. Then say how the caller
knows it lost.
\end{bankq}

\begin{bankq}{This query is slow. What do you do?}
\qasked
Usually a follow-up once a data model is on the board.

\qtested
Whether you have an order, and whether you know what an index costs.

\qstaff
Look at the plan before changing anything \dash{} the fix depends on whether it
is a scan, a bad join order, or a correct plan over too much data.

Then: an index on the columns actually filtered, in an order that matches how
they are filtered; covering the selected columns where that removes a lookup;
and an awareness that every index is paid for on every write, so an index
nobody's query uses is a pure cost.

\judgement{The answer that reads as senior is the last sentence, because it
treats an index as a trade rather than a fix, and because it implies you have
been on the wrong side of it.}

\qsenior
``Add an index'', correct in most cases, offered without looking at the plan or
mentioning the write cost.

\qcontested
Little at this level.

\qdrill
For a query you know, say which index you would add and name the query that
would then be slower.
\end{bankq}

\begin{bankq}{How would you change this API without breaking its callers?}
\qasked
Asked of any interface with consumers you do not control, which at this level
is most of them.

\qtested
Whether you know expand-and-contract, and whether you know who controls the
timing.

\qstaff
Expand, migrate, contract. Add the new form alongside the old; write to both or
translate between them; move callers; only then remove the old one \dash{} and
the removal happens when the last caller has moved, which is a fact you need
evidence for rather than a date you choose.

The part that separates levels is the last clause. You do not control your
callers' release cycles, so ``we will remove it in the next version'' is a plan
that depends on other people's calendars. Instrumenting the old path to find
out whether anyone still uses it is the difference between a migration and a
hope.

\qsenior
Versioning the endpoint. It works and it defers the problem: you now maintain
two implementations and the question of when the old one dies is unanswered.

\qcontested
Whether to version at all, and where. A real and long-running disagreement.

\qdrill
Say how you would know that nobody is calling the old form any more.
\end{bankq}

\begin{bankq}{What breaks first if this gets ten times the traffic?}
\qasked
Nearly universal in a design round, and the same probe as the scale question
in the project deep-dive.

\qtested
Whether you have an order of attack. This is the same question as
Chapter~\ref{ch:deep-dive}'s probe about scale, and it is asked in almost every
design round.

\qstaff
Answer in order rather than listing. The database connection pool is usually
first and is usually the cheapest to fix. Then the database itself \dash{} and
which part: write throughput, read throughput, or lock contention on one row,
which are three different problems with three different fixes. Then anything
that holds state in a process, which is the point at which horizontal scaling
stops being free. Then any external dependency with its own limit, where your
capacity is somebody else's rate limit. Then, usually last, the thing everyone
names first, which is the application's own processing.

\judgement{Naming the order, and saying what evidence would change it, is the
answer. Listing the components is not, and the difference is audible.}

\qsenior
A correct list of components that could become bottlenecks, unordered.

\qcontested
Nothing general \dash{} the order is a property of the system, and saying so is
part of the answer.

\qdrill
For a system you have operated, say what actually broke first the last time
load went up. It is frequently not what you would have predicted, and that
story is worth more in an interview than the general answer.
\end{bankq}

\section*{What this chapter does not claim}

No claim is made about how often these appear relative to the domain-specific
surfaces. The reason they are here is structural: a build round involves data,
concurrency and an interface, so these questions are reachable from any
exercise, and the observation that candidates under-prepare them is the
author's.
